\documentclass{article}
\usepackage{amsmath, amssymb, amsthm}
\usepackage[a4paper,
            bindingoffset=0.2in,
            left=1in,
            right=1in,
            top=1in,
            bottom=1in,
            footskip=.25in]{geometry}
\newtheorem{theorem}{Teorema}[section]
\renewcommand\qedsymbol{$\blacksquare$}
\renewenvironment{proof}{{\bfseries Dem:}}

\title{Ejercicios}
\author{Josè Hernàndez}
\date{\today}

\begin{document}
\maketitle

\section*{Practica 4}

\subsection*{Ejercicio 1} \hfill

Sea \( f \colon X \to \mathbb{R} \) y \( \mathcal{C} \subset \mathcal{B}(\mathbb{R}) \) una clase de subconjuntos
que genera a \( \mathcal{B}(\mathbb{R}) \). Demuestre que \( f \) es medible si, y solo si,
\( f^{-1}(\mathcal{C}) \subseteq \mathcal{A} \).
\bigskip

\begin{proof}

   \( \left( \implies \right)  \) \\
   Sea \( f \) medible. Sabemos que, para todo \( A \in \mathcal{B}(\mathbb{R}) \),
   \( f^{-1}(A) \in \mathcal{A} \).\\
   Sea \( B \in f^{-1}(\mathcal{C}) \), entonces existe \( C \in \mathcal{C} \) tal que 
   \( B = f^{-1}(C) \), pero como \( \mathcal{C} \subseteq \mathcal{B}(\mathbb{R}) \), \( C \in \mathcal{B}(\mathbb{R}) \)
   es decir, \( B = f^{-1}(C) \in \mathcal{A} \), para todo \( B \in f^{-1}(\mathcal{C}) \), asì, 
   \( f^{-1}(\mathcal{C}) \subseteq \mathcal{A}\). 

   \( \left( \impliedby \right) \) \\
   Supongamos que \( f^{-1}(\mathcal{C}) \subseteq \mathcal{A} \), queremos ver que 
   \( \forall B \in \mathcal{B}(\mathbb{R}) \), \( f^{-1}(B) \in \mathcal{A} \) (\( f \) es medible).\\
   Sea \( M = \{B \in \mathcal{B}(\mathbb{R}) \mid f^{-1}(B) \in \mathcal{A}\}   \), es claro que
   \( M \subseteq \mathcal{B}(\mathbb{R}) \). Notemos que \( \mathcal{C} \subseteq M \) por
   consecuencia de la hipotesis,
   si \( M \) es una \( \sigma -algebra \) entonces, \( \mathcal{B}(\mathbb{R}) = \sigma (\mathcal{C}) \subseteq M \),
   haciendo que \( M = \mathcal{B}(\mathbb{R}) \), obteniendo el resultado. Veamos que \( M \) es una \( \sigma -algebra \)

    1) Notemos que \( f^{-1}(\mathbb{R}) = X \in \mathcal{A} \)

    2) Sea \( A \in M \), \( f^{-1}(A^{c}) = (f^{-1}(A))^{c} \in \mathcal{A} \), y como 
    \( A^{c} \in \mathcal{B}(\mathbb{R}) \), \( A^{c} \in M \)

    3) Sea \( \{A_n\} \subseteq M \), queremos ver que \( \bigcup_{n = 1}^{\infty}{A_n} \in M \).
    \( f^{-1}(\bigcup_{n = 1}^{\infty}{A_n}) = \bigcup_{n = 1}^{\infty}{f^{-1}(A_n)} \in \mathcal{A} \),
    y como \(  \bigcup_{n = 1}^{\infty}{A_n} \in \mathcal{B}(\mathbb{R})\), \( \bigcup_{n = 1}^{\infty}{A_n} \in M\)

    Asi, hemos visto que \( M \) es una \( \sigma -algebra \), por lo que se obtiene el resultado buscado.
    \( \hfill \blacksquare \)

\end{proof}





\clearpage

\subsection*{Ejercicio 2}
\bigskip

\begin{proof}

    (a) Sea \( B \in \mathcal{B}(\mathbb{R}) \),

    \begin{equation*}
        h^{-1}(B) = 
        \begin{cases}
            h^{-1}(\{\alpha \} ) & \text{si \( \alpha \in B \)} \\
            \emptyset & \text{si no}
        \end{cases}
    \end{equation*}

    pero, como \( h^{-1}(\{\alpha \}  ) = X \), se cumple que, en cada caso,
    \( h^{-1}(B) \in \mathcal{A} \) 
    \bigskip

    (b) Sea \( a \in \mathbb{R} \), usando el Ejercicio 1, sabemos que 
    \( \mathcal{C} = \{(-\infty, a) \mid a \in \mathbb{R}\} \) es una clase generadora de \( \mathcal{B}(\mathbb{R}) \),
    entonces \( f+g \) sera borel-medible sii, \( (f+g)^{-1}(\mathcal{C}) \in \mathcal{A} \). Veamos lo sig.
    \begin{align*}
        (f+g)^{-1}(-\infty,a) &= \{x \in X \mid f(x)+g(x) < a \}  \\
                              &= \{x \in X \mid f(x) < a-g(x)\}  
    \end{align*}
    \begin{align*}
        f(x) < a - g(x) &\iff \exists q \in \mathbb{Q} \text{ tq } f(x) < q < a - g(x) \\
                        &\iff f(x) < q \text{ y } q < a - g(x) \\
                        &\iff \exists q \in \mathbb{Q} \text{ tq } x \in \{x \in X \mid  f(x) < q\} \cap \{x \in X \mid q <a -g(x)\}    
    \end{align*}

    De donde,

    \begin{align*}
        (f+g)^{-1}(-\infty,a) &= \bigcup_{q \in \mathbb{Q}} \{x \in X \mid  f(x) < q\} \cap \{x \in X \mid q <a -g(x)\} \\
                              &= \bigcup_{q \in \mathbb{Q}} f^{-1}(-\infty,q) \cap g^{-1}(-\infty, a-q)
    \end{align*}

    Luego, como \( f \) y \( g \) son medibles, \( f^{-1}(-\infty,q) , g^{-1}(-\infty, a-q) \in \mathcal{A}\), por lo tanto, como \( \mathbb{Q} \) es numerable,
    \(  (f+g)^{-1}(-\infty,a) \in \mathcal{A}\), asi, \( f+g \) es medible.
    \bigskip

    (c) Por el mismo argumento que en la parte (b), veamos \( (\alpha f^{-1})(-\infty,a) \), \( a \in \mathbb{R} \).
    Tenemos dos casos,

    1) Si \( \alpha = 0 \), \( \alpha f = 0 \), \(  (\alpha f^{-1})(-\infty,a) = \{x \in X \mid 0 < a\}  \)
    \begin{align*}
        (\alpha f^{-1})(-\infty,a) &= \{x \in X \mid 0 < a\} \\ 
                                   &= 
                                   \begin{cases}
                                       X & \text{si } a > 0 \\
                                       \O & \text{si } a \le 0
                                   \end{cases}
    \end{align*}
    En cualquier caso, \( (\alpha f^{-1})(-\infty,a) \in \mathcal{A} \)

    2) Si \( \alpha \neq 0 \),
    \begin{align*}
        (\alpha f^{-1})(-\infty,a) &= \{x \in X \mid \alpha f(x) < a\} \\
                                   &= \{x \in X \mid  f(x) < \frac{a}{\alpha}\} \\
                                   &= f^{-1}(-\infty, \frac{a}{\alpha})  
    \end{align*}
    Luego, como \( f \) es medible, se concluye que \( (\alpha f^{-1})(-\infty,a) \in \mathcal{A} \), y asi,
    \( \alpha f \) es medible.
    \bigskip

    (g) Notemos que 
    \begin{align*}
        (max \{f,g\}^{-1}) (-\infty,a) &= \{x \in X \mid f(x) < a\} \cup \{x \in X \mid g(x) < a\} \\
                                       &= f^{-1}(-\infty,a) \cup g^{-1}(-\infty,a) \in \mathcal{A}
    \end{align*}

    De donde, \( max \{f,g\} \) es medible.

    Es analogo para \( min \{f,g\} \), ya que \( \mathcal{A} \) es cerrado bajo intersecciones numerables.
    \bigskip

    (h) Podemos definir \( f^+ \) como \( max \{f,0\} \), de donde, usando el ejercicio (g), se obtiene el resultado.
    
    Analogamente para \( f^- = min \{f,0\}   \).


\end{proof}

\clearpage

\section*{Práctica 6}

\begin{itemize}
    \item[2.] Demuestre que \( \bigotimes_{i=1}^{n} \mathcal{B}(\mathbb{R}) = \mathcal{B}(\mathbb{R}^n) \)

        \begin{proof}

            Basta ver que ambas \( \sigma -algebras \) son el resultado de la \( \sigma -algebra \)
            generada por la misma clase de conjuntos \( \mathcal{C} \). O en su defecto, ver que coinciden en
            al menos un conjunto generador.
            
            El conjunto generador de \( \bigotimes_{i=1}^{n} \mathcal{B}(\mathbb{R}) \) es la clase de los bloques medibles.

            Falta analizar \( \mathcal{B}(\mathbb{R}^n) \)
        \end{proof}

\end{itemize}

\end{document}

